\chapter{Reference: lifecycle hooks}
\label{ch:refman-lifecycle}

\begin{aside}
Three moments in a program's life are special: start, every
update, and exit. \maya{} exposes them as commands and as
the natural shape of \code{init} and \code{update}.
\end{aside}

\section{Where each hook lives}

\begin{itemize}[leftmargin=*]
  \item \textbf{Start} — \code{init()} returns the initial
    Model. Pair it with a startup command by returning from
    \code{init} into the loop via a follow-up message.
  \item \textbf{Every update} — \code{update()} runs on
    every message; that's your "tick" if you want one.
  \item \textbf{Exit} — return \code{Cmd<Msg>::quit()}; the
    runtime calls destructors as \code{main} unwinds.
\end{itemize}

\section{Startup command pattern}

If you need to fire a command at startup, define a
\code{Boot} variant and dispatch it as the first message:

\begin{lstlisting}[language=C++]
#include <maya/maya.hpp>
#include <chrono>

using namespace maya;
using namespace maya::dsl;

struct Boot {
    struct Model { std::string status = "booting..."; };
    struct Start {};
    struct Ready {};
    struct Quit {};
    using Msg = std::variant<Start, Ready, Quit>;

    static Model init() { return {}; }

    static auto update(Model m, Msg msg)
        -> std::pair<Model, Cmd<Msg>>
    {
        return std::visit(overload{
            [&](Start) {
                return std::pair{m,
                    Cmd<Msg>::after(std::chrono::seconds(1),
                                    Ready{})};
            },
            [&](Ready) {
                m.status = "ready";
                return std::pair{m, Cmd<Msg>{}};
            },
            [](Quit) {
                return std::pair{Model{}, Cmd<Msg>::quit()};
            },
        }, msg);
    }

    static Element view(const Model& m) {
        return v(
            text(m.status) | Bold,
            blank_,
            t<"q to quit"> | Dim
        ) | pad<1> | border_<Round>;
    }

    static auto subscribe(const Model&) -> Sub<Msg> {
        return Sub<Msg>::batch(
            key_map<Msg>({ {'q', Quit{}} }),
            // Send Start once on first frame.
            Sub<Msg>::on_startup(Start{})
        );
    }
};

int main() { run<Boot>({.title = "boot"}); }
\end{lstlisting}

\section{Notes}

\begin{itemize}[leftmargin=*]
  \item Use \code{on\_startup} for one-shot init tasks;
    after the first frame, it's a no-op.
  \item Cleanup belongs in destructors. Don't try to do
    "exit work" inside the \code{Quit} handler beyond
    setting flags.
\end{itemize}

\section{Recap}

\begin{recap}
\begin{itemize}[leftmargin=*]
  \item \code{init} for startup state.
  \item \code{on\_startup} for one-shot commands.
  \item RAII for exit work.
\end{itemize}
\end{recap}

\section*{Workshop}

\begin{enumerate}[leftmargin=*]
  \item Load a config file on startup, show errors inline.
  \item Print a goodbye card after \code{run} returns.
  \item Persist Model state to disk on quit.
\end{enumerate}
